\begin{abstract}
Registramos el electrocardiograma de un voluntario con un módulo AD8232 y un Arduino UNO, durante \SI{1}{\minute} en reposo y \SI{3}{\minute} inmediatamente después de una actividad física intensa, y analizamos la variabilidad de su frecuencia cardíaca. Determinamos a partir del firmware y de los registros que la frecuencia de muestreo es de \SI{500}{\hertz}. Acondicionamos la señal con filtros de fase cero, excluimos los tramos saturados y la discontinuidad del inicio, detectamos los latidos con una versión de Pan-Tompkins de umbral fijo y estimamos el espectro del tacograma con el periodograma de Lomb. Después del ejercicio la frecuencia cardíaca media pasó de 84 a \SI{131}{lpm} y bajó de 146 a \SI{120}{lpm} durante la recuperación. El RMSSD cayó de 48 a \SI{8.6}{\milli\second} y las potencias LF y HF se redujeron entre 50 y 100 veces. La banda HF casi desapareció en el primer minuto y volvió a aparecer en los siguientes. Estos resultados coinciden con la retirada vagal y su reactivación gradual descritas en la bibliografía, con las limitaciones de un único voluntario y de registros cortos para estimar la banda LF.
\end{abstract}
